\begin{alist}
\item Η γραφική παράσταση της $f$ διέρχεται από το $A(-1,3)$ και δεν θα μπορούσε να διέρχεται και από το σημείο $B(2,5)$, διότι η $f$ έχει πεδίο ορισμού και σύνολο τιμών το $\mathbb{R}$, είναι όμως γνησίως φθίνουσα και θα έπρεπε να ισχύει $-1 < 2 \Rightarrow f(-1) > f(2)$. Όμως $f(-1)=3$ και $f(2)=5$.

\item Η γραφική παράσταση της $f$ θα μπορούσε να είναι η ii. διότι ενώ όλες διέρχονται από το σημείο $(-1,3)$, η ii. είναι γραφική παράσταση γνησίως φθίνουσας συνάρτησης.
\end{alist}
